\section{Chương~\ref{sec:sensors}. Đầu vào của cỗ máy}

Cấu tạo của cảm biến đã được đo trực tiếp: bằng bản ghi dòng điện của từng tế bào, dao động của ốc tai và việc đếm tế bào võng mạc; giới hạn độ nhạy và công thức nhiễu hạt suy ra từ vật lý, còn việc mã của cảm biến được tinh chỉnh để mang nhiều thông tin nhất là giả thuyết có cơ sở vững.

{\footnotesize
\setlength{\tabcolsep}{3pt}\renewcommand{\arraystretch}{1.15}
\begin{longtable}{>{\raggedright}p{0.5cm}>{\raggedright}p{4.0cm}>{\raggedright}p{5.6cm}>{\raggedright}p{2.3cm}>{\raggedright\arraybackslash}p{1.7cm}}
\toprule
TT & Khẳng định & Thí nghiệm và điều đã đo & Nguồn & Trạng thái \\
\midrule\endfirsthead
\toprule
TT & Khẳng định & Thí nghiệm và điều đã đo & Nguồn & Trạng thái \\
\midrule\endhead
1 & Cảm biến là bộ chuyển đổi: thụ thể tạo ra dòng điện, còn đầu ra của tế bào cảm giác ở mắt và tai là glutamate đưa lên bus \nt{GLUT}; các tế bào đầu tiên không phát xung (hình~\ref{fig:transducer}) & Tế bào que và tế bào nón của rùa giải phóng một amino acid kích thích: nó được các kênh glutamate trên một mảnh màng tách rời bắt được. Tế bào lông trong của ốc tai chuột cống: dòng điện ở đầu tận cùng sợi thần kinh đi qua thụ thể glutamate AMPA, tần số giải phóng tăng theo khử cực liên tục của tế bào tới $150$\,Hz & \cite{CopenhagenJahr1989,GlowatzkiFuchs2002} & đã đo \\
2 & Trong bóng tối, cảm biến ánh sáng đáp ứng với một photon bằng dòng khoảng một picoampe & Điện cực hút trên đoạn ngoài tế bào que của cóc: đáp ứng với chớp sáng mờ bị lượng tử hóa, mỗi sự kiện $\approx1$\,pA ($5\%$ dòng tối) với độ tản $0.2$\,pA, hiệu suất $0.5$. Người báo nhận được một photon đúng nhiều hơn mức ngẫu nhiên & \cite{Baylor1979,Tinsley2016} & đã đo \\
3 & Cảm biến âm thanh nhận ra độ dịch chuyển nhỏ hơn một nanomét & Phương pháp Mössbauer, màng đáy ốc tai chuột lang tại vị trí $16$--$19$\,kHz: ở ngưỡng đáp ứng của dây thần kinh thính giác, vận tốc màng khoảng $0.04$\,mm/s, tức độ dịch chuyển $\approx0.35$\,nm (chúng tôi quy đổi). Đo trên màng, không phải trên bó lông của cảm biến & \cite{Sellick1982,Hudspeth1989} & đã đo \\
4 & Trong bóng tối, độ chính xác bị giới hạn bởi nhiễu hạt photon~\eqref{eq:photonnoise}; ánh sáng yếu thì tích lũy lâu hơn & Công thức suy ra từ thống kê Poisson. Ở tế bào que, nhiễu dưới ánh sáng mờ trùng với tổng các sự kiện lượng tử độc lập; trên nền sáng, đáp ứng với chớp sáng nhỏ hơn và ngắn hơn. Con số ``vài phần mười giây'' ở đây chưa được kiểm tra riêng & \cite{Baylor1979} & lý thuyết \\
5 & Dải đầu vào là mười bậc với ánh sáng và mười hai bậc với âm thanh, còn tần số xung chưa tới ba bậc & Với âm thanh: đáp ứng của màng đáy ở tần số đặc trưng tăng với độ dốc tới $0.2$\,dB/dB trong dải $40$--$80$\,dB, sự nén vẫn thấy được trên $100$\,dB. Bản thân số bậc là ước tính chuẩn, trong chương không dẫn nguồn & \cite{Ruggero1997} & đã đo \\
6 & Đáp ứng của cảm biến theo~\eqref{eq:nakarushton}, còn $\sigma$ theo độ sáng trung bình, làm đường cong dịch dọc trục logarit (hình~\ref{fig:adaptation}) & Công thức lần đầu được khớp với điện thế S của võng mạc cá. Tế bào nón của rùa: đáp ứng tuân theo $V=I V_m/(I+\sigma)$ và bao $\sim3.5$ bậc độ sáng; sau $2$--$3$ phút có nền, đường cong dịch ngang và giữ nguyên độ rộng. Liên hệ với phép chia cho hoạt động chung: tổng quan & \cite{NakaRushton1966,NormannPerlman1979,Carandini2012} & đã đo \\
7 & Cảm biến truyền các thay đổi, và mã của chúng được tinh chỉnh để mang nhiều thông tin nhất trên mỗi xung (mệnh đề~\ref{prop:sensors}) & Nguyên lý được đề xuất về lý thuyết. Bằng chứng mạnh nhất: ở ruồi, đặc tuyến ``độ tương phản $\to$ đáp ứng'' của nơ-ron lớp thứ nhất trùng với phân phối tích lũy độ tương phản của cảnh tự nhiên, và nhờ đó đạt thông tin lớn nhất & \cite{Barlow1961,Laughlin1981} & giả thuyết \\
8 & Cảm biến bật và cảm biến tắt là mã hai dây; camera sự kiện lặp lại nó trên silicon & Tổng quan thí nghiệm: kênh bật và kênh tắt của võng mạc tách nhau ngay từ khớp thần kinh đầu tiên và có thể bị vô hiệu hóa riêng rẽ. Camera: $128\times128$ điểm ảnh không có khung hình, dải $120$\,dB, độ trễ $15$\,\textmu s & \cite{Schiller1992,Lichtsteiner2008,Gallego2022} & đã đo \\
9 & Mắt có $\sim10^8$ cảm biến ánh sáng và $\sim10^6$ nơ-ron đầu ra: nén $\sim100$ lần & Đếm trên toàn bộ tiêu bản võng mạc người: trung bình $92$ triệu tế bào que và $4.6$ triệu tế bào nón; tế bào đầu ra (tế bào hạch) từ $0.7$ đến $1.5$ triệu tùy mắt & \cite{Curcio1990,CurcioAllen1990} & đã đo \\
10 & Ở chuột nhắt, mắt có hơn ba mươi kênh đầu ra & Chuột nhắt: chụp calci hai photon hơn $11\,000$ tế bào đầu ra và phân cụm đáp ứng: rõ ràng nhiều hơn $30$ loại chức năng. Ở người con số này chưa được đo & \cite{Baden2016} & đã đo \\
11 & Mắt chuột lang truyền $\sim875\,000$ bit/s; quy đổi cho mắt người được $\sim10^7$ bit/s & Chuột lang, mảng đa điện cực, chuyển động trong cảnh tự nhiên: $6$--$13$ bit/s mỗi tế bào, $\sim875\,000$ bit/s cho $\sim10^5$ tế bào đầu ra. Với người ($\sim10^6$ tế bào), $\sim10^7$ bit/s là quy đổi của các tác giả & \cite{Koch2006} & đã đo \\
12 & Tai là dãy bộ lọc thông dải với bản đồ tần số gần như logarit (hình~\ref{fig:filterbank}) & Vị trí dao động lớn nhất của màng đáy phụ thuộc tần số; hàm ``tần số--vị trí'' gần như hàm mũ, cùng một dạng mô tả ốc tai người và động vật sống & \cite{Greenwood1990,Robles2001} & đã đo \\
13 & Bộ cộng hưởng là chủ động: một phần cảm biến khuếch đại dao động yếu hàng nghìn lần về biên độ ($66$--$76$\,dB), âm to lên thì hệ số khuếch đại giảm & Đo rung bằng laser ở đáy ốc tai sóc chinchilla: với âm yếu ở tần số đặc trưng, hệ số khuếch đại $66$--$76$\,dB so với xương bàn đạp; sau khi con vật chết, độ nhạy giảm $60$--$81$\,dB ($10^3$--$10^4$ lần về biên độ) và sự nén biến mất. Nguồn năng lượng là tế bào lông ngoài & \cite{Ruggero1997,Robles2001} & đã đo \\
14 & Bộ khuếch đại của tai đặt sát ranh giới tự kích & Tính toán: gần điểm rẽ nhánh Hopf, không tránh khỏi nén dải, chọn lọc tần số sắc nét và âm tổ hợp, tức mọi phi tuyến đã biết của thính giác. Việc ốc tai được chỉnh đúng như vậy chưa được chứng minh trực tiếp & \cite{Eguiluz2000} & giả thuyết \\
15 & Ở tần số thấp, xung đi cùng pha với âm (ở chuột lang, khóa pha yếu đi từ $\sim600$\,Hz và còn thấy tới $\sim3.5$\,kHz), và từ đó tìm được hướng; ở tần số cao chỉ còn mã vị trí & Dây thần kinh thính giác chuột lang: khóa pha bắt đầu giảm từ khoảng $600$\,Hz và biến mất trên $3.5$\,kHz; ở mèo và khỉ, ngưỡng này cao hơn khoảng một quãng tám. Chênh lệch thời điểm âm đến hai tai: tổng quan & \cite{PalmerRussell1986,Grothe2010} & đã đo \\
16 & Các cảm biến khác (bảng~\ref{tab:sensors}): khứu giác có $\sim400$ loại thụ thể, mỗi nơ-ron một loại, và mã tổ hợp; vị giác đi một phần qua ATP và serotonin; xúc giác có cảm biến thích nghi nhanh và chậm; nóng và lạnh tách riêng & Chuột nhắt: một thụ thể nhận nhiều mùi, một mùi kích hoạt nhiều thụ thể, các mùi khác nhau cho các tập khác nhau. Không có thụ thể ATP thì dây thần kinh vị giác không đáp ứng; nụ vị giác giải phóng serotonin. Ghi từng sợi thần kinh ở da bàn tay người: bốn loại theo tốc độ thích nghi. Kênh lạnh khác với các kênh nóng & \cite{Mombaerts2004,Malnic1999,Finger2005,Huang2005,JohanssonVallbo1979,McKemy2002} & đã đo \\
\bottomrule
\end{longtable}}
